\section{Extending \DVRFOne and \glitter to be Robust}
\label{robust-shine}

As presented, \DVRFOne is not a robust algorithm: if corrupt players
output incorrect values from $\DVRFOnem.\generate$, then
$\DVRFOnem.\verify$ will compute that the $|\coalition| \ge 2t-1$
commitments do not all lie on the same degree-$(t-1)$ polynomial in the
exponent, and output 0.  However, $\DVRFOnem.\verify$ will not
necessarily be able to identify \emph{which} party or parties
misbehaved.  This is because it is possible for the $t-1$ corrupt
parties to output commitments that are wrong, but consistent with up to
$t-1$ of the honest players' commitments, making it look like the
remaining honest player was the outlier.  Therefore, as we see in
Figure~\ref{fg:construction},
\glitter will have to just abort in the case that
$\DVRFOnem.\verify$ outputs 0.

However, we observe that in the \emph{honest supermajority} case, where
$|\coalition| \ge 3t-2$, and so there are at most $t-1$ corrupt parties
and at least $2t-1$ honest parties, $\DVRFOnem.\verify$ can be modified
to be robust: if any corrupt parties output incorrect values from
$\DVRFOnem.\generate$, $\DVRFOnem.\verify$ will be able to notice the
inconsistency \emph{and identify the misbehaving parties}.  Those
parties can then be kicked out of $\coalition$, and \glitter will be
able to continue and produce its correct deterministic signature.  The
misbehaving parties may also be able to be kicked out of the system
entirely, or have other external sanctions applied.
The same observation was made by Cramer et al.~\cite{CramerDI05} for their NIVSS, but in their setting, the \emph{shares} are themselves published, whereas in \DVRFOne, only \emph{commitments} are available to $\DVRFOnem.\verify$.

The reason this works is based on a simple fact about polynomials, as
used in error-correcting codes: there can only be at most one polynomial
of degree $t-1$ that passes through at least $2t-1$ of a given set of
$3t-2$ points.  The reason is that if $f_1$ and $f_2$ each passed
through a (possibly different) set of at least $2t-1$ of the $3t-2$
points, then those two sets must intersect in at least $t$ points (by
simple counting).  Those $t$ points uniquely define a polynomial of
degree $t-1$, and so $f_1 = f_2$.  In the honest supermajority setting,
the at least $2t-1$ honest parties will give consistent outputs from
$\DVRFOnem.\generate$, and so there will be exactly one (not at most
one) polynomial of degree $t-1$ that passes through at least $2t-1$ of
the exponents of the $3t-2$ commitments.

The typical way to \emph{find} this unique polynomial is with the
Berlekamp-Welch algorithm, but $\DVRFOnem.\verify$ only has access to
commitments to the purported polynomial evaluations, and not the
evaluations themselves.  However, we can take advantage of the fact that
we are already assuming that ${n-1 \choose t-1}$ is not too large, and
take a more brute force approach.  The modification to
$\DVRFOnem.\verify$ is then as follows: if player $\partyid$ running
$\DVRFOnem.\verify$ would output 0 in Step 4 of $\DVRFOnem.\verify$, instead iterate
through all $t-1$-sized subsets $S$ of $\coalition \setminus
\{\partyid\}$, and reconstruct (in the exponent) the polynomial that
passes through player $\partyid$'s own point, and the points of $S$.  If
that polynomial fails to pass through at most $t-1$ of the points, those
points are the misbehaving parties.  Exclude them from $\coalition$, and
output 1.

There are ${|\coalition|-1 \choose t-1}$ possible subsets $S$, and at
least ${2t-1 \choose t-1}$ of those subsets contain only honest parties.
If the iteration is done in a pseudorandom order (for example, keyed by
$\vrfprivout_\partyid$), then the expected number of iterations before the misbehaving
parties are identified is at most ${|\coalition|-1 \choose t-1} / {2t-1
\choose t-1}$, which is significantly smaller than ${n-1 \choose t-1}$,
which we already assumed is not too large.

%\chelesa{I don't want to make this claim.}
%In addition, potentially malicious parties are disincentivized from
%giving incorrect outputs from $\DVRFOnem.\generate$ exactly
%\emph{because} they will be able to be identified and sanctioned, and so
%even a moderately large (but feasible) computation cost to identify
%misbehavers may be acceptable.

Extending \DVRFOne to be robust will make round 1 of \glitter robust, so
that the correct $R$ can be computed even if some or all of the
malicious parties submit incorrect $R_j$ values.
To make round 2 of \glitter also robust, we must also guard against an
adversary submitting a correct $R_j$ in round 1, but an incorrect
signature share $z_j$ in round 2.

In the event of the failure of the test $g^z \neq R \cdot \pk^c$ in
$\combine$, one can identify the parties that submitted incorrect
$z_j$ values by checking $g^{z_i} \iseq R_i \cdot (\pk_i)^c$.  Once the
parties submitting incorrect signature shares are eliminated, the
remaining correct shares (of which there will be at least $t$) can be
combined into the final signature $\sigma=(R,z)$.

\subsection{Communication Model for Robust \glitter}

For non-robust \glitter,
we discuss in Remark~\ref{remark:auth-channel} how our adversarial model requires authenticated channels,
but does not assume replay protection.
However,
additionally in the robust setting,
parties will require a communication channel that ensures protection against replayed messages,
to ensure that an adversary cannot replay old messages from honest participants,
resulting in a potential abort.
Furthermore, to achieve robustness, the communication model likewise requires liveness;
the adversary cannot prevent honest parties from exchanging messages.

